% E3: Broader Impact Statement
% Addresses societal implications of automated LLM routing.

\subsection{Broader Impact}
\label{appendix:broader_impact}

\paragraph{Cost reduction and accessibility.}
Adaptive routing reduces the average cost of LLM inference by directing
easy prompts to cheaper models.  Our multi-model experiments quantify this
potential: at $\lambda{=}1.0$ with a $K{=}3$ portfolio, banditGPT achieves
substantial cost reduction over static frontier-model routing
\textbf{TODO: update with K=3 cost reduction numbers}.  If widely deployed, this could lower the cost barrier
for organizations and individuals who currently cannot afford frontier
models for all queries, thereby broadening access to high-quality AI
assistance.

\paragraph{Energy efficiency.}
By routing fewer prompts to large, compute-intensive models, banditGPT
reduces the aggregate energy consumption and carbon footprint of LLM serving.
At scale, even small reductions in the fraction of queries sent to frontier
models can translate into meaningful resource savings.

\paragraph{Quality-of-service risks.}
Automated routing introduces a potential failure mode: the router may
systematically under-route certain prompt types to the stronger model.
If prompt difficulty correlates with user demographics or domains (e.g.,
non-English prompts, specialized technical queries), this could create
disparate quality of service.  The Corralling safety mechanism
($\gamma$-mixing floor) ensures all models continue to receive a minimum
fraction of traffic, preventing complete abandonment of any model.
Additionally, the calibration analysis (Appendix~\ref{appendix:calibration})
shows the router's predictions are well-calibrated (bias $< 0.05$),
reducing systematic misrouting risk.

Multi-provider portfolios provide an additional resilience benefit: at
$K{=}3$, multi-provider routing distributes traffic across providers, so a single-provider
outage affects only a fraction of requests---a property absent from static
single-model deployments.
Practitioners should nonetheless audit routing patterns for demographic
fairness before production deployment.

\paragraph{Judge model dependency.}
All routing decisions ultimately depend on the LLM judge (GPT-4/GPT-4o) that
produces the training signal.  Biases in the judge---preference for verbose
outputs, position bias, cultural biases---propagate through the reward signal
into the router's learned policy.  banditGPT is designed to be modular with
respect to the reward source: the reward signal can be replaced with human
feedback, A/B test outcomes, or an ensemble of judges without changing the
routing algorithm.  This modularity allows practitioners to select or combine
reward sources that best align with their quality objectives.

\paragraph{Dual use.}
The bandit-based routing framework is model-agnostic and could be applied to
route queries between any set of models, including models with different
safety profiles.  A malicious operator could configure the router to maximize
a reward signal that conflicts with user safety (e.g., routing sensitive
queries to an unaligned model).  This risk is not specific to our work but is
inherent to any programmable routing layer.  The system's transparency
properties (below) partially mitigate this by making routing decisions
auditable.

\paragraph{Transparency.}
The contextual bandit framework provides interpretability advantages over
black-box routing: the LinUCB parameters ($\hat{\theta}_a$) reveal which
embedding directions drive routing decisions, and the Corralling weights
indicate which expert strategy the system trusts.  This enables
administrators to inspect and audit routing policies.  However, the
sentence embeddings themselves are opaque, making it difficult to explain
to end users \emph{why} a particular model was chosen for their query.
Future work could address this through attention-based embedding methods
or post-hoc feature attribution.
